% =============================================================
%  American Terawatt — SMP Level 1, single-line diagram
%
%  Structure:
%    smp-level-1.tex             <- this file
%    theme -> ../theme           <- symlink: theme/atstyle.sty and theme/fonts/
%                                   resolve as in the main presentation
%    .latexmkrc                  <- adds ~/hids/template/glyphwright to TEXINPUTS
%    figures/figures.yaml        <- figure list; figures/smp_level_1.tex is
%    figures/src/smp_level_1.yaml   generated from the diagram YAML
%
%  Compile in this folder:
%    latexmk -xelatex smp-level-1.tex
%
%  The symbols come from the Glyphwright library in ~/hids/template/glyphwright
%  (\usepackage{glyphwright}); regenerate figure and library as described in
%  figures/figures.yaml.
% =============================================================
\documentclass[aspectratio=169,11pt]{beamer}
\usepackage{theme/atstyle}
\usepackage{glyphwright}

\title{HVDC Design}
\subtitle{AmWatt Team Presenation}
\author{American Terawatt}
\date{\today}

\newcommand{\slfig}[1]{\input{figures/#1}}

\begin{document}

\titleslide

% ---- Single-line diagram ------------------------------------
\begin{frame}{System Overview}
  ORC generators feed a data centre through a point-to-point VSC HVDC link.
  \vspace{2.2em}

  \centering
  \slfig{smp_level_1}
\end{frame}

% ---- Power path ---------------------------------------------
\begin{frame}{SMP Level 1\\Power path}
  \begin{columns}[T,onlytextwidth]
    \begin{column}{0.50\textwidth}
      \begin{block}{From generation to load}
        A.c.\ collection, conversion to d.c., d.c.\ transmission, conversion back to a.c.
      \end{block}
      \vspace{0.3em}
      \begin{itemize}
        \setlength{\itemsep}{0.4em}
        \item Generation: three ORC sets on one busbar.
        \item Sending end: transformer YNd11, VSC converter.
        \item D.c.\ link: two poles, cable sealing ends at both ends.
        \item Receiving end: VSC converter, transformer, data centre.
      \end{itemize}
    \end{column}
    \begin{column}{0.44\textwidth}
      \footnotesize
      \renewcommand{\arraystretch}{1.1}
      \begin{tabularx}{\linewidth}{@{}l >{\raggedright\arraybackslash}X@{}}
        \toprule
        \multicolumn{1}{c}{\textbf{Tag}} & \textbf{Element} \\
        \midrule
        GEN1--3 & ORC generator set \\
        BUS1    & Busbar, sending end \\
        T1      & Transformer YNd11 \\
        U1      & VSC converter, sending end \\
        CSE1--4 & Cable sealing ends, d.c.\ link \\
        U2      & VSC converter, receiving end \\
        T2      & Transformer YNd11 \\
        BUS2    & Busbar, receiving end \\
        DATA    & Data centre \\
        \bottomrule
      \end{tabularx}
    \end{column}
  \end{columns}
\end{frame}

\end{document}
